\section{Open Challenges and Future Directions}
\label{sec:challenges}

% === Consolidated Discussion by Tier ===

\subsection{External Controls}

Context engineering, guardrails, and decoding strategies share a common strength: they treat the model as a black box, making them model-agnostic, composable, and rapidly updatable without retraining. Context engineering is training-free and flexible, with different techniques naturally addressing distinct challenges---prompting for lightweight alignment, RAG for factual grounding, memory for persistence, and reasoning for self-regulation. Guardrails add an independent policy-enforcement layer that can be swapped across providers, while decoding strategies enable fine-grained, real-time steering of token distributions across multiple trustworthiness dimensions.

However, all three tiers are fundamentally soft controls and share similar failure modes. Context engineering is fragile to small phrasing changes and vulnerable to context poisoning; guardrails face a continuous cat-and-mouse dynamic against adversarial evasion; and hard decoding constraints can produce brittle, incoherent outputs~\citep{Huang2024SafeAlignerSA}. A common ``Control Tax'' underlies these methods: strict guardrails increase false positives and reduce helpfulness, while aggressive logit suppression~\citep{Zhao2024AdversarialCD} or noise injection~\citep{Majmudar2022DifferentiallyPD} compromises generation quality. Complex LLM-based judges and multi-layer decoding further introduce significant latency overhead. None of these methods can provide formal safety guarantees in isolation, and integrated pipelines that combine context engineering, guardrails, and decoding remain underexplored.

\subsection{Internal Manipulations}

Representation engineering, unlearning, and pruning offer more direct behavioral control by intervening on internal model components. Representation engineering treats hidden states as controllable objects, enabling concept-level steering through activation addition or sparse autoencoders~\citep{zou2023representation,li2023iti,turner2023steering}. Unlearning steers outputs away from sensitive knowledge through runtime penalties or influence subtraction, making it practical for dynamic regulatory demands. Pruning provides structural control by removing heads, neurons, or tokens associated with unsafe or biased behavior, and is compatible with broader compression pipelines.

Despite these strengths, internal methods share key fragilities. Representation engineering depends heavily on intervention strength, layer choice, and normalization; multi-concept steering remains brittle, and strong interventions risk pushing activations out of distribution~\citep{wehner2025taxonomy}. Unlearning faces a fundamental verification problem: many methods suppress rather than truly remove knowledge, leaving models vulnerable to adversarial recovery~\citep{lucki2024adversarial}. Pruning is highly sensitive to where sparsity is introduced---even small amounts of naive pruning can sharply weaken safety-critical behavior. All three methods also exhibit a utility--control trade-off, where stronger interventions can degrade fluency, accuracy, or reasoning~\citep{liu2024breaking}. Additionally, representation engineering carries dual-use risks, as the same activation-level controls can strengthen or weaken safety~\citep{revisiting2025jailbreaking}. Robust deployment across internal methods will require better verification protocols, principled intervention selection, and careful integration with broader safety governance~\citep{nist2025ai8001}.

\subsection{System-Level Orchestration}

Multi-agent systems introduce reliability through collaborative reasoning, but their most immediate obstacle is operational overhead: each query may trigger multiple agents over several rounds, with performance gains often diminishing quickly~\citep{conf/coling/0008L0HF25, 10.5555/3692070.3692537}, rendering many architectures impractical for real-time applications~\citep{song2024auditllmmultiagentcollaborationlogbased}. Moreover, the collaborative dynamics designed to foster trustworthiness can become vectors for emergent failures. Agents sharing the same base model exhibit ``groupthink,'' reinforcing shared hallucinations~\citep{conf/coling/0008L0HF25}, and adversaries can exploit the collaborative structure---a single persuasive adversary can sway an entire group~\citep{amayuelas-etal-2024-multiagent}, while debate frameworks can be more vulnerable to jailbreaks than single agents~\citep{qi2025amplifiedvulnerabilitiesstructuredjailbreak}. Evaluating MAS trustworthiness is also fundamentally more complex, as effectiveness depends on roles, topology, protocols, and agent count~\citep{zhu2025mastermultiagentsecurityexploration}.


%%%%%%%%%%%%%%%%%%% Original Disucssion (DO NOT TOUCH)%%%%%%%%%%%%%%%%%%%
% \textcolor{blue}{Below are original Discussion Section}
% The inference-time methods surveyed in this work represent a paradigm shift towards dynamic, modular, and accessible control over LLM behavior. While promising, their widespread adoption is hindered by fundamental challenges that cut across the entire taxonomy. Moving from ad-hoc interventions to a robust engineering discipline for trustworthiness requires a concerted research effort focused on overcoming the fragility of controls, managing performance trade-offs, developing scalable solutions, and understanding system-level effects. We synthesize these challenges into a forward-looking research agenda.

% \paragraph{Overcoming the Brittleness of Controls}
% A recurring theme is the fragility of inference-time interventions. Controls are often highly sensitive to minor variations in input or model state, leading to unpredictable failures. Context engineering, for instance, is notoriously fragile, where slight rephrasing can bypass safety instructions. This brittleness extends to internal methods; knowledge editing has revealed that LLMs are highly fragile to weight perturbations, where localized edits can cause catastrophic degradation of general abilities. Similarly, representation engineering struggles with fragile generalization, as steering effects can be inconsistent across different layers or generation lengths, while hard constraints in decoding strategies can lead to brittle and incoherent outputs. Future work must move beyond controls that work only under narrow conditions. This necessitates developing methods for \textit{robustification}, potentially through adversarial training of the control mechanisms themselves, to ensure they generalize across a wider distribution of inputs and contexts.

% \paragraph{Quantifying and Mitigating the Control Tax}
% Nearly every method introduces a trade-off between trustworthiness and utility, a phenomenon that can be termed a "Control Tax." Imposing stricter controls on model behavior often degrades core performance metrics like fluency, accuracy, and helpfulness. This is explicit in guardrails, which must balance safety and helpfulness, and in decoding strategies, which face a direct trade-off between control strength and generation quality. The tension is also evident in unlearning methods, where achieving stronger unlearning efficacy often harms model utility, and in pruning, where removing components for safety can cause collateral damage to general capabilities. A primary goal for future research should be to formally understand and quantify this Control Tax. This requires creating standardized benchmarks that explicitly measure performance across a joint distribution of task-based and trust-based metrics, with the ultimate aim of designing new methods that push the Pareto frontier, minimizing the utility penalty required to achieve a desired level of trust.

% \paragraph{Addressing the Scalability and Efficiency Barrier}
% For inference-time methods to be practical, they must be computationally efficient. Several powerful techniques, however, impose significant overhead in latency and cost, creating a "viability barrier" for real-world deployment. Multi-agent systems are a primary example, as their multiplicative demand for LLM calls renders them impractical for most real-time applications. Similarly, complex guardrails, particularly those using an LLM-as-a-judge, introduce noticeable latency that is unacceptable in many user-facing scenarios. Even advanced decoding strategies can introduce computational overhead at each generation step. The field therefore needs a greater focus on lightweight and efficient control mechanisms. Promising directions include distilling the behavior of complex MAS or judge models into smaller, faster classifiers and developing pruning techniques that are aligned with hardware acceleration to achieve speedups without compromising trustworthiness.

% \paragraph{Defending Against Adaptive Adversaries}
% Many current methods are static defenses vulnerable to adaptive adversaries who continuously discover new bypass techniques. This creates a reactive "cat-and-mouse" dynamic, where developers are constantly patching vulnerabilities. Guardrails are often susceptible to adversarial evasion, and unlearning methods can be circumvented by adversarial probing that recovers supposedly erased information. More complex systems are not immune; multi-agent systems can be manipulated by a single persuasive adversary and have been shown to be more vulnerable to certain jailbreaks than single models. Future work must therefore shift from static defenses to dynamic, adaptive security postures. This involves developing methods for real-time threat detection, automated red-teaming to discover vulnerabilities proactively, and control systems that can adapt their strategies in response to detected attacks.

% \paragraph{The Unsolved Problems of Generalization and Composition}
% A critical limitation of many targeted interventions is their failure to generalize. Furthermore, how different inference-time methods interact when layered together remains poorly understood. Knowledge editing exemplifies the generalization challenge, where an edit often fails to transfer to rephrased queries and simultaneously disrupts unrelated knowledge, demonstrating poor locality. Representation engineering faces similar issues, where a steering vector may not generalize across contexts or may interfere with other concepts due to entangled representations. This leads to the broader challenge of \textit{compositionality}. The current ad-hoc approach of layering controls risks interventions canceling each other out or creating unexpected failure modes. A priority should be developing a principled science of composition, including theoretical frameworks to predict interactions and evaluation suites that specifically test the robustness of combined systems.

% \paragraph{Toward Principled Evaluation and Formal Guarantees}
% The evaluation of trustworthiness remains a significant bottleneck. The development of benchmarks like HalluEditBench has shown that the real-world efficacy of techniques such as knowledge editing can be far lower than suggested by earlier, simpler datasets. Moreover, most methods lack formal guarantees about their behavior. For instance, unlearning and pruning cannot currently guarantee that information has been truly erased rather than merely suppressed. The field urgently needs more robust, comprehensive, and standardized evaluation protocols that test for long-horizon consistency, adversarial robustness, and complex multi-objective trade-offs. A critical long-term goal is to move from heuristic controls to methods with provable guarantees, particularly for non-negotiable domains like privacy and safety.

% \paragraph{Understanding and Mitigating Emergent System-Level Risks}
% Finally, intervening in a complex system like an LLM can have unintended, systemic consequences. The collaborative dynamics of advanced methods can introduce novel and dangerous failure modes. Multi-agent systems, for example, are susceptible to emergent risks like "groupthink," where shared biases are amplified, and "cascading errors," where one agent’s mistake misleads the entire system. A similar systemic risk is seen in sequential knowledge editing, where a series of seemingly localized edits can accumulate distortions that lead to a sudden collapse of the model's general abilities. A systems-thinking approach is therefore essential. Instead of analyzing interventions in isolation, future research must model and test for emergent behaviors and feedback loops, with the ultimate goal of building not just trustworthy components, but robustly trustworthy systems.
